% GENERATED FILE. Edit canonical lessons, then run npm run generate:books.
% canonical-source-hash: fnv1a64:a39f076358d622ef
% canonical-lessons: UR-C177-mulazim, UR-C177-sathi, UR-C177-mazdur, UR-C177-kisan, UR-C177-darzi

\chapter{Employee, Colleague, Labourer, Farmer, Tailor}
\label{ch:ur-employee-colleague-labourer-farmer-tailor}
\hlchaptermodality{177}

\begin{chapteropening}
\textbf{By the end of this chapter:} \emph{I can say an employee, a colleague, a labourer, a farmer and a tailor.}

Complete the last lesson of chapter 177: I can say an employee, a colleague, a labourer, a farmer and a tailor.
\end{chapteropening}

\section[\texorpdfstring{\ur{ملازم} (mulāzim)}{ملازم (mulāzim)}]{\ur{ملازم} (mulāzim) — an employee}

\label{lesson:UR-C177-mulazim}



\begin{quote}
Before the new one: say the Urdu for a rupee, then the Urdu for a coin.
\end{quote}

\subsection*{You'll want to know: \ur{ملازم}}
\textbf{\ur{ملازم}} — \emph{mulāzim} — \textquotedblleft{}an employee\textquotedblright{}.

\begin{cousinweb}
\textbf{\ur{ملازم}} (\emph{mulāzim}) is a loan from Arabic \emph{mulāzim}, \textquotedblleft{}one who keeps close, an attendant\textquotedblright{}, which reached Urdu through Persian. Its Arabic root is \emph{l-z-m} (Lane).
\end{cousinweb}

\begin{grammarlens}[title={What you've built}]
One more word for this chapter.
\end{grammarlens}

\subsection*{Your turn}
\begin{itemize}
\raggedright
  \item \emph{Say it:} \emph{mulāzim}
  \item \emph{Say it:} \emph{mulāzim}, once more
  \item \emph{Say it:} say \emph{mulāzim}
  \item \emph{Recall:} say the Urdu for a rupee, then the Urdu for a coin, then say \emph{mulāzim} again
\end{itemize}

\begin{tcolorbox}[breakable,colback=teal!4,colframe=teal!35!black,title={Before you move on}]
What does \ur{ملازم} mean? (\textquotedblleft{}An employee\textquotedblright{}.) Say it once more.
\end{tcolorbox}
\section[\texorpdfstring{\ur{ساتھی} (sāthī)}{ساتھی (sāthī)}]{\ur{ساتھی} (sāthī) — a colleague, a companion}

\label{lesson:UR-C177-sathi}



\begin{quote}
Before the new one: say the Urdu for a coin, then the Urdu for an employee.
\end{quote}

\subsection*{You'll want to know: \ur{ساتھی}}
\textbf{\ur{ساتھی}} — \emph{sāthī} — \textquotedblleft{}a colleague, a companion\textquotedblright{}.

\begin{cousinweb}
\textbf{\ur{ساتھی}} (\emph{sāthī}) is a tadbhava, a word inherited and reshaped by sound change. It comes from \emph{*sārthin}, \textquotedblleft{}having part in a company\textquotedblright{}, a form CDIAL reconstructs from its descendants (CDIAL 13366). It is built on \emph{sārtha}, \textquotedblleft{}company\textquotedblright{}.
\end{cousinweb}

\begin{grammarlens}[title={What you've built}]
One more word for this chapter.
\end{grammarlens}

\subsection*{Your turn}
\begin{itemize}
\raggedright
  \item \emph{Say it:} \emph{sāthī}
  \item \emph{Say it:} \emph{sāthī}, once more
  \item \emph{Say it:} say \emph{sāthī}
  \item \emph{Recall:} say the Urdu for a coin, then the Urdu for an employee, then say \emph{sāthī} again
\end{itemize}

\begin{tcolorbox}[breakable,colback=teal!4,colframe=teal!35!black,title={Before you move on}]
What does \ur{ساتھی} mean? (\textquotedblleft{}A colleague, a companion\textquotedblright{}.) Say it once more.
\end{tcolorbox}
\section[\texorpdfstring{\ur{مزدور} (mazdūr)}{مزدور (mazdūr)}]{\ur{مزدور} (mazdūr) — a labourer}

\label{lesson:UR-C177-mazdur}



\begin{quote}
Before the new one: say the Urdu for an employee, then the Urdu for a colleague.
\end{quote}

\subsection*{You'll want to know: \ur{مزدور}}
\textbf{\ur{مزدور}} — \emph{mazdūr} — \textquotedblleft{}a labourer\textquotedblright{}.

\begin{cousinweb}
\textbf{\ur{مزدور}} (\emph{mazdūr}) is a loan, usually traced to Persian \emph{muzdūr}, \textquotedblleft{}a hired labourer\textquotedblright{}. It is built on \emph{muzd}, \textquotedblleft{}wages\textquotedblright{}.
\end{cousinweb}

\begin{grammarlens}[title={What you've built}]
One more word for this chapter.
\end{grammarlens}

\subsection*{Your turn}
\begin{itemize}
\raggedright
  \item \emph{Say it:} \emph{mazdūr}
  \item \emph{Say it:} \emph{mazdūr}, once more
  \item \emph{Say it:} say \emph{mazdūr}
  \item \emph{Recall:} say the Urdu for an employee, then the Urdu for a colleague, then say \emph{mazdūr} again
\end{itemize}

\begin{tcolorbox}[breakable,colback=teal!4,colframe=teal!35!black,title={Before you move on}]
What does \ur{مزدور} mean? (\textquotedblleft{}A labourer\textquotedblright{}.) Say it once more.
\end{tcolorbox}
\section[\texorpdfstring{\ur{کسان} (kisān)}{کسان (kisān)}]{\ur{کسان} (kisān) — a farmer}

\label{lesson:UR-C177-kisan}



\begin{quote}
Before the new one: say the Urdu for a colleague, then the Urdu for a labourer.
\end{quote}

\subsection*{You'll want to know: \ur{کسان}}
\textbf{\ur{کسان}} — \emph{kisān} — \textquotedblleft{}a farmer\textquotedblright{}.

\begin{cousinweb}
\textbf{\ur{کسان}} (\emph{kisān}) is a tadbhava, a word inherited from Sanskrit and reshaped by sound change. It comes from \emph{kṛṣāṇa}, \textquotedblleft{}ploughing\textquotedblright{} (CDIAL 3447). It is the present participle of \emph{kṛṣati}, \textquotedblleft{}ploughs\textquotedblright{}: a farmer is \textquotedblleft{}the one ploughing\textquotedblright{}.
\end{cousinweb}

\begin{grammarlens}[title={What you've built}]
One more word for this chapter.
\end{grammarlens}

\subsection*{Your turn}
\begin{itemize}
\raggedright
  \item \emph{Say it:} \emph{kisān}
  \item \emph{Say it:} \emph{kisān}, once more
  \item \emph{Say it:} say \emph{kisān}
  \item \emph{Recall:} say the Urdu for a colleague, then the Urdu for a labourer, then say \emph{kisān} again
\end{itemize}

\begin{tcolorbox}[breakable,colback=teal!4,colframe=teal!35!black,title={Before you move on}]
What does \ur{کسان} mean? (\textquotedblleft{}A farmer\textquotedblright{}.) Say it once more.
\end{tcolorbox}
\section[\texorpdfstring{\ur{درزی} (darzī)}{درزی (darzī)}]{\ur{درزی} (darzī) — a tailor}

\label{lesson:UR-C177-darzi}



\begin{quote}
Before the new one: say the Urdu for a labourer, then the Urdu for a farmer.
\end{quote}

\subsection*{You'll want to know: \ur{درزی}}
\textbf{\ur{درزی}} — \emph{darzī} — \textquotedblleft{}a tailor\textquotedblright{}.

\begin{cousinweb}
\textbf{\ur{درزی}} (\emph{darzī}) is a loan, usually traced to Persian \emph{darzī}, \textquotedblleft{}tailor\textquotedblright{}. It is built on \emph{darz}, \textquotedblleft{}seam\textquotedblright{}.
\end{cousinweb}

\begin{grammarlens}[title={What you've built}]
One more word for this chapter.
\end{grammarlens}

\subsection*{Your turn}
\begin{itemize}
\raggedright
  \item \emph{Say it:} \emph{darzī}
  \item \emph{Say it:} \emph{darzī}, once more
  \item \emph{Say it:} say \emph{darzī}
  \item \emph{Recall:} say the Urdu for a labourer, then the Urdu for a farmer, then say \emph{darzī} again
\end{itemize}

\begin{tcolorbox}[breakable,colback=teal!4,colframe=teal!35!black,title={Before you move on}]
What does \ur{درزی} mean? (\textquotedblleft{}A tailor\textquotedblright{}.) Say it once more.
\end{tcolorbox}
